\section{The characteristic-\texorpdfstring{$p$}{p} argument}\label{sec:charp}

In characteristic $p$, Frobenius makes the differential entries of
the complex from Proposition~\ref{prop:setup} arbitrarily deep.
We apply Theorem~\ref{thm:complex} separately to each component of
$T$ and add the resulting inequalities. The following
Hilbert--Kunz identity identifies that sum with $e(S)$.

\begin{lemma}[Ma's Hilbert--Kunz identity]\label{lem:ma-hk}
Let $(T,\mathfrak t)$ be a complete local ring of characteristic
$p>0$ with perfect residue field, let $J$ be a perfect ideal, and put
$S=T/J$.  For any system of parameters $x_1,\ldots,x_d$ of $S$,
\[
 e_{\mathrm{HK}}(J+(x),T)
   =\sum_{\substack{Q\text{ minimal over }J\\\dim(T/Q)=d}}
        e((x),T/Q)\,e_{\mathrm{HK}}(JT_Q,T_Q).
\]
If each $T_Q$ occurring in this sum is a complete intersection, then
\[
       e_{\mathrm{HK}}(J+(x),T)=e((x),S).
\]
\end{lemma}

\begin{proof}
The first assertion is \cite[Lemma~4.2]{Ma}. For the second, we use
the following complete-intersection consequence from
\cite[Section~2.3 and Lemma~4.2]{Ma}: a primary ideal of finite
projective dimension in a local complete intersection of
characteristic $p$ has Hilbert--Kunz multiplicity equal to its
colength. Since $J$ is perfect, $JT_Q$ is such an ideal in $T_Q$.
Consequently
\[
 e_{\mathrm{HK}}(JT_Q,T_Q)=\length_{T_Q}(T_Q/JT_Q).
\]
The Hilbert--Samuel associativity formula gives the conclusion.  No
assumption that $x$ extend to a parameter system of $T$ is involved.
\end{proof}

\begin{theorem}\label{thm:positive-characteristic}
For every flat local map $(R,\mathfrak m)\to(S,\mathfrak n)$ of
nonzero Noetherian local rings of characteristic $p>0$, one has
$e(R)\leq e(S)$.
\end{theorem}

\begin{proof}
Lemma~\ref{lem:reduction} reduces the problem to
complete rings of the same positive dimension, with $R$ a domain and
the residue field of $S$ algebraically closed.  Use
Proposition~\ref{prop:setup} to obtain $T,J,P_\bullet,F_\bullet$ and
the integers $d,b,h=d+b$.  Set
\[
       H=J+(x_1,\ldots,x_d)T.
\]
This is a $\mathfrak t$-primary ideal.

Fix a minimal prime $P$ of $T$ and put $D=T/P$, with maximal ideal
$\mathfrak d$.  This is a complete local domain of dimension $h$ and
has the same algebraically closed residue field.  Choose a minimal
reduction $(z_1,\ldots,z_h)$ of $\mathfrak d$.  There is an integer
$c\geq0$ such that
\[
  \mathfrak d^n=(z)^{n-c}\mathfrak d^c
      \subseteq(z)^{n-c}\qquad(n\geq c).
\]
Take the $D$-algebra $C$ and its normalized length $\lambda$ from
Theorem~\ref{thm:length-inputs}.  For $I=(z)C$ the length hypotheses
of Theorem~\ref{thm:complex} hold with
\[
       \mu=e((z),D)=e(D).
\]

For a power $q=p^n$, let $F_\bullet(q)$ be the complex obtained by
raising every matrix entry of the differentials of $F_\bullet$ to
its $q$th power.  This is scalar extension by the $n$th Frobenius
endomorphism of $T$ and has the same ranks as $F_\bullet$.  Its
differential entries lie in $\mathfrak t^{[q]}$.  Their images in
$C$ therefore lie in
\[
     \mathfrak d^{[q]}C\subseteq\mathfrak d^qC
                          \subseteq I^{q-c}.
\]
Frobenius preserves the identities defining a contracting homotopy
on each localization where $F_\bullet$ is contractible.  Subsequent
scalar extension preserves those identities as well.  Since $(z)$
is $\mathfrak d$-primary, the ideals $I$ and $\mathfrak tC$ have
the same radical.  Consequently
$F_\bullet(q)\otimes_T C$ is contractible locally off $V(I)$.

The first differential of the resolution $P_\bullet$ has entries
generating $J$, and $P_0=T$.  Thus
\[
 H_0(F_\bullet(q)\otimes_T C)=C/H^{[q]}C.
\]
Apply Theorem~\ref{thm:complex} with $s=q-c$, for $q>c$.  Its
degree-zero rank is one, and the ranks of all its terms are fixed
as $q$ varies.  The Hilbert--Kunz quotient formula from
Theorem~\ref{thm:length-inputs} and Frobenius scaling give
\begin{align*}
 q^h e_{\mathrm{HK}}(HD,D)
   &=e_{\mathrm{HK}}((HD)^{[q]},D)\\
   &=\lambda(C/H^{[q]}C)\\
   &\geq e(D)(q-c)^h(1-o(1)).
\end{align*}
Dividing by $q^h$ and letting $q$ tend to infinity proves
\[
             e_{\mathrm{HK}}(HD,D)\geq e(D).
\]

Every component of $T$ has dimension $h$ by
Proposition~\ref{prop:setup}.  The associativity formulas for
Hilbert--Kunz~\cite[Theorem~3.14]{HunekeHK} and Hilbert--Samuel
multiplicities have the same
generic length coefficients, so
\begin{align*}
 e_{\mathrm{HK}}(H,T)
   &=\sum_{P\in\operatorname{Min}(T)}
       \length_{T_P}(T_P)\,
       e_{\mathrm{HK}}(H(T/P),T/P)\\
   &\geq\sum_{P\in\operatorname{Min}(T)}
       \length_{T_P}(T_P)\,e(T/P)
     =e(T).
\end{align*}
In particular this summation does not assume that $T$ is reduced.

We now verify the hypotheses of Lemma~\ref{lem:ma-hk}.  The ring $T$
is complete of characteristic $p$, its residue field is algebraically
closed and therefore perfect, and $J$ is perfect.  For every minimal
prime $Q$ of $J$ with $\dim(T/Q)=d$, Proposition~\ref{prop:setup}
shows that $T_Q$ is a complete intersection.  The chosen $x$ is a
system of parameters of $S=T/J$, and its ideal is a reduction of
$\mathfrak n$.  Hence
\[
  e(S)=e((x),S)=e_{\mathrm{HK}}(H,T)
          \geq e(T)=e(R),
\]
as required.
\end{proof}

\subsection{Dutta multiplicity of short complexes}

The zeroth homology in the following consequence need not be cyclic.
The finite-length-module comparison in
Lemma~\ref{lem:perfection-finite-length} therefore takes the place of
the Hilbert--Kunz quotient formula used above.
We use the following terminology of Iyengar--Ma--Walker
\cite[Sections~2.5--2.6]{IMW2022}.

\begin{definition}\label{def:short-complex}
Let $(R,\mathfrak m)$ be a Noetherian local ring of characteristic
$p>0$ and dimension $d$.  A \emph{short complex} over $R$ is a
complex of finite free modules
\[
 F=(0\longrightarrow F_d\longrightarrow\cdots
                 \longrightarrow F_0\longrightarrow0)
\]
whose homology modules have finite length and for which $H_0(F)\ne0$.
For any bounded complex $G$ with finite-length homology, write
\[
 \chi_R(G)=\sum_i(-1)^i\length_R H_i(G).
\]
When $G$ is finite free, let $\Phi_R^nG$ denote scalar extension
along the $n$th Frobenius endomorphism.  Its \emph{Dutta multiplicity}
is
\[
 \chi_\infty^R(G)=
 \lim_{n\longrightarrow\infty}
       p^{-nd}\chi_R(\Phi_R^nG).
\]
\end{definition}

The limit exists; see \cite[Section~4.17]{IMW2022}.  For a short
complex over a complete characteristic-$p$ local ring, the Frobenius
homology theorem of Roberts, as recalled in
\cite[Theorem~2.4]{Ma}, states that
\[
 \lim_{n\longrightarrow\infty}
 p^{-nd}\length_R H_i(\Phi_R^nF)=0\qquad(i>0).
\]
Consequently
\begin{equation}\label{eq:dutta-h0}
 \chi_\infty^R(F)=
 \lim_{n\longrightarrow\infty}
 p^{-nd}\length_R H_0(\Phi_R^nF).
\end{equation}
This holds for arbitrary short complexes; it does not require their
zeroth homology to be cyclic.

Iyengar, Ma, and Walker conjectured the lower bound in the next
corollary for short complexes over all complete local rings
\cite[Conjecture~8.1]{IMW2022}.  Here we prove the bound for complete
domains of positive characteristic.

\begin{corollary}[Dutta multiplicity over positive-characteristic domains]
\label{cor:dutta-domain}
Let $(R,\mathfrak m)$ be a complete Noetherian local domain of
characteristic $p>0$, and let $F$ be a short complex over $R$.
Then
\[
                    \chi_\infty^R(F)\geq e(R).
\]
\end{corollary}

\begin{proof}
Put $d=\dim R$.  If $d=0$, then $R$ is a field and $F_0$ is a
nonzero free module, so $\chi_\infty^R(F)=\rank_R F_0\geq1=e(R)$.
Assume $d>0$.  Splitting off contractible free summands makes $F$
minimal without changing its homology or Dutta multiplicity.
Write $b_i=\rank_R F_i$; in particular $b_0\geq1$.

First suppose that the residue field is algebraically closed.  Choose
a minimal reduction $I=(z_1,\ldots,z_d)$ of $\mathfrak m$.  As in the
proof of Theorem~\ref{thm:positive-characteristic}, there is a fixed
integer $c\geq0$ with
\[
                 \mathfrak m^q\subseteq I^{q-c}\qquad(q\geq c).
\]
Use the characteristic-$p$ algebra $C=R_{\mathrm{perf}}$ and normalized
length $\lambda$ from Theorem~\ref{thm:length-inputs}.  For any short
complex $G$ over $R$, right exactness gives
\[
 C\otimes_R H_0(G)=H_0(C\otimes_R G),\qquad
 \Phi_R^n(H_0(G))=H_0(\Phi_R^nG).
\]
Lemma~\ref{lem:perfection-finite-length} and
\eqref{eq:dutta-h0} therefore give
\[
                  \lambda H_0(C\otimes_R G)=\chi_\infty^R(G).
\]

Let $q=p^n$.  The entries of the differentials of $\Phi_R^nF$ lie
in $\mathfrak m^{[q]}\subseteq\mathfrak m^q\subseteq I^{q-c}$.
Frobenius preserves a contracting homotopy at every localization
where $F$ is contractible.  Thus $\Phi_R^nF$ again has finite-length
homology; it is minimal and has nonzero zeroth homology, so it is
short.  Scalar extension to $C$ preserves the same contracting
homotopies; since $I$ is $\mathfrak m$-primary, the extended complex
is contractible locally off $V(IC)$.  The defining Dutta limit, with
its index shifted by $n$, now gives
\[
 \lambda H_0(C\otimes_R\Phi_R^nF)
   =\chi_\infty^R(\Phi_R^nF)=q^d\chi_\infty^R(F).
\]
Moreover $\sum_i(-1)^ib_i=0$, since $F$ becomes exact over the
fraction field of $R$.  Apply Theorem~\ref{thm:complex} to
$C\otimes_R\Phi_R^nF$ with parameter ideal $IC$, depth $s=q-c$,
and $\mu=e(I,R)=e(R)$.  For $q>c$ it gives
\[
 q^d\chi_\infty^R(F)
 \geq e(R)(q-c)^d
       \left(b_0-\frac{A_d}{q-c}\sum_{i=0}^d b_i\right).
\]
After division by $q^d$, the right side tends to $b_0e(R)$.
Since $b_0\geq1$, this proves the corollary in this case.

For an arbitrary residue field, Lemma~\ref{lem:residue-extension}
gives a complete flat local extension
$(R,\mathfrak m)\to(S,\mathfrak n)$ with
$\mathfrak n=\mathfrak mS$, algebraically closed residue field, and
$\dim S=d$, $e(S)=e(R)$.  Put $G=S\otimes_R F$.  This is again a
minimal short complex.  Frobenius scalar extension commutes with
this base change, and finite-length module lengths are preserved,
so
\[
             \chi_\infty^S(G)=\chi_\infty^R(F);
\]
this invariance is also recorded in \cite[Section~4.17]{IMW2022}.
The ring $S$ need not be a domain, so we pass to its components by
a prime filtration rather than apply the preceding case directly.

Choose a finite prime filtration of the $S$-module $S$ with factors
$D_j=S/P_j$, and put $h_j=\dim D_j$.  Additivity of Euler
characteristic for the bounded free complex $\Phi_S^nG$ gives
\[
 \chi_S(\Phi_S^nG)
       =\sum_j\chi_S(D_j\otimes_S\Phi_S^nG).
\]
The complexes in the sum have finite-length homology because the
contracting homotopies off $\mathfrak n$ persist after tensoring.
The Hochster--Huneke Frobenius homology estimate, in the form
\cite[Theorem~5.2]{BhattHochsterMa}, gives constants $C_j$ such that
\[
 \length_S H_i(D_j\otimes_S\Phi_S^nG)
       \leq C_jp^{n\min\{h_j,d-i\}}
       \qquad(0\leq i\leq d).
\]
Indeed, its complex has length $d$ and finite-length homology, and
its finite coefficient module $D_j$ has dimension $h_j$.
When $h_j<d$, every homology length in this summand is
$O(p^{nh_j})$, so its contribution divided by $p^{nd}$ tends to
zero.

When $h_j=d$, the ring $D_j$ is a complete local domain with the
same algebraically closed residue field.  The complex
$D_j\otimes_SG$ is short over $D_j$: its homology is supported at
the maximal ideal because the localized contracting homotopies
persist, and its zeroth homology is nonzero because its matrices
remain minimal and $b_0\geq1$.  Inspecting differential matrices
gives
\[
 D_j\otimes_S\Phi_S^nG
       \cong\Phi_{D_j}^n(D_j\otimes_SG).
\]
Lengths of its homology over $S$ and $D_j$ agree.  Taking limits,
applying the already proved domain case to these factors, and using
additivity of top-dimensional Hilbert--Samuel multiplicity along the
prime filtration, we obtain
\[
 \chi_\infty^R(F)=\chi_\infty^S(G)
   =\sum_{h_j=d}\chi_\infty^{D_j}(D_j\otimes_SG)
   \geq\sum_{h_j=d}e(D_j)=e(S)=e(R).
\]
This completes the proof.
\end{proof}
